\chapter{Comparative Analysis}
\label{06_Comparative_Analysis}
\thispagestyle{empty}

With all models created, this chapter focuses on their comparative analysis to highlight performance across different dimensions. This evaluation is structured into five distinct sections, each focusing on a critical aspect of the model's behavior.

\section{Training Efficiency and Computational Costs}

In the context of training our classifiers on our extensive dataset, one of the most crucial differences is training complexity and computational cost. In order to reach the best metrics, it is required to use as many training entries as possible and also fine-tune large space of hyperparameters. This, in itself, makes the process very challenging, even with strong Google Colaboratory resources. These challenges are particularly important when training deep learning models (CNN, BiLSTM, BERT).

\vspace{3mm}

In contrast, classic machine learning models like Logistic Regression and SVM have simpler architectures and consequently are less resource-intensive. This simplicity is reflected in the size of the final models: the Logistic Regression model is a lightweight 4.4 MB, while the CNN model occupies over 717 MB. Significantly larger size of deep learning models underscores their ability to capture complex patterns but also highlights challenges in storage and deployment.

\vspace{3mm}

GPU acceleration is essential for training deep learning models efficiently. Unlike CPUs, most commonly limited up to 20 cores, GPUs can perform thousands of parallel operations, dramatically reducing training time. Classic models relied on local CPU training, as they did not fully benefit from GPU capabilities. This made them faster to train on a CPU, but less suitable for leveraging the advantages of GPUs, especially with larger~datasets. 

\vspace{3mm}

In summary, GPU usage is essential for efficiently training modern neural networks. In machine learning, the ability to leverage GPU computation is a critical factor. Without GPUs, training models such as BERT would either take an impractical amount of time or require the computational power of supercomputers. This reliance highlights the critical role of GPUs in advancing AI research, enabling scientists to build more sophisticated models at a larger scale.

\section{Text Representations}

Bag-of-Words (BoW) and Term Frequency-Inverse Document Frequency (TF-IDF) are foundational text representation methods that identify textual similarity based on shared vocabulary. However, these approaches do not account for word order, semantic relationships, or word location, which are crucial for deep textual understanding. For instance, sentences like "I love dogs" and "Dogs love me" would generate almost identical vector representations in a BoW model, despite their different meanings.

\vspace{3mm}

While TF-IDF shares this limitation, it adjusts word frequencies within a document to reflect how often a term appears relative to the document's length. Additionally, it assigns weights based on the term's rarity across the dataset, giving higher importance to words that are specific to a subset of documents. This dual adjustment makes TF-IDF more effective for ranking terms by relevance in tasks like information retrieval.


\vspace{3mm} % !!!!

Even though BoW and TF-IDF are effective for basic text representation and analysis, their highlighted limitations leave room for more advanced approaches. Modern natural language processing tasks increasingly rely on word embeddings, which address these issues by capturing relationships between words in a continuous vector space. This advancement represents an important shift from frequency-based representations to context-aware vectorization techniques.

\vspace{3mm}

Word embeddings, such as those generated by fastText, are a powerful approach to capturing the semantic relationships between words, sentences, or entire documents. Unlike traditional text representation methods, such as BoW and TF-IDF which treat words as discrete, independent entities, word embeddings represent words as dense vectors in a continuous vector space. These vectors encode rich semantic and syntactic information, allowing for better relationship retrieval.

\vspace{3mm}

A classic example of this can be seen in Figure~\ref{fig:word_embeddings}, which illustrates the analogy "king is to queen as man is to woman." In the figure, words like "king", "queen", "man" and "woman" are represented as vectors in a continuous vector space. These vectors are not just positioned close to each other due to semantic similarity but also maintain parallel relationships, showing their analogous meaning.

\begin{figure}[h!]
    \centering
    \includegraphics[width=\textwidth]{Images/word_embeddings.png}
    \caption{Visualization of word embeddings, including the analogy "king is to queen as man is to woman" There are 4 analogies one can construct, based on the parallel red arrows and their
direction. \cite{sutor2019metaconcepts}.}
    \label{fig:word_embeddings}
\end{figure}

For instance, the vector difference between "king" and "man" is approximately the same as the vector difference between "queen" and "woman." This relational similarity is indicated by the parallel red arrows in the figure, demonstrating that the analogy "king + woman - man = queen" holds true. This kind of relationship enables word embeddings to detect analogies that traditional models like BoW or TF-IDF cannot perform. Even if the vectors are not perfectly aligned, the embeddings are designed to reflect the closest and most meaningful relationships in the data.

\vspace{3mm}

This ability to encode complex relationships in a high-dimensional space highlights the strength of word embeddings in capturing nuances and contextual meanings in text, making them much more effective than traditional methods.

\vspace{3mm}

In our case, word embeddings generated by fastText significantly enhanced the performance of BiLSTM and CNN models by providing a more meaningful representation of words that is aware of the context in the text. Unlike BoW or TF-IDF, embeddings are capable of capturing contextual nuances, such as the difference in meaning between "bank" in "river bank" and "bank" in "financial institution", depending on the surrounding words.

\vspace{3mm}

Additionally, fastText stands out among word embedding techniques by modeling both entire words and their subword structures. Unlike Word2Vec and GloVe, which fail to handle out-of-vocabulary (OOV) words, fastText generates embeddings using character-level n-grams. This allows the approximation of representations for unseen words. This capability is more robust and ensures greater generalization, making it particularly effective for datasets with diverse vocabularies. 

\vspace{3mm}

Lastly, BERT uses contextualized word embeddings that address the limitations of earlier methods. Unlike static embeddings, where a word always has the same single vector, BERT dynamically generates embeddings based on the surrounding words. For example, the word "bank" would have two different embeddings in "river bank" and "financial bank", making classification less reliant on further model architecture. BERT’s bidirectional approach further enhances its understanding by simultaneously considering both preceding and succeeding words.

\vspace{3mm}

Moreover, the fine-tuning capability of BERT allows it to excel in a wide variety of tasks, from sentiment analysis and question answering to classification. It represents the current state-of-the-art for text representation in NLP, setting new benchmarks for performance. However, these advantages come at a certain cost: BERT is significantly more computationally intensive than simpler methods like fastText or TF-IDF. In our case, even the considerable computational resources of Google Colaboratory were insufficient for extensive experimentation with its architecture. Thus, while BERT offers superior performance, selecting a text representation method should carefully take into account these trade-offs.

\section{Model Interpretations and Context Extraction}

\vspace{3mm}

The representation of words through methods like TF-IDF and Bag-of-Words is particularly effective for Logistic Regression and SVM models when there are big disparities in word frequency. These methods leverage linear and hyperplane-based decision-making processes of these models. 

\vspace{3mm}

To analyze this, we compute the frequency distribution of words across human and AI-generated text subsets, normalizing for differences in the average number of words per text. From these distributions, we extracted the 20 most frequent words from the human subset and matched their corresponding frequencies in the AI subset.

\vspace{3mm}


Figure \ref{fig:most_popular_words} illustrates these findings. Most frequent words contribute a lot to the overall text composition, with their cumulative frequency summing to 27.29\% in human texts and 25.40\% in AI texts. Each AI bar is annotated with the relative frequency of the word compared to the human subset.

\vspace{3mm}

\begin{figure}[h!]
    \centering
    \includegraphics[width=\textwidth]{Images/most_popular_words.png}
    \caption{Comparison of Word Frequencies in Human and AI-Generated Texts.}
    \label{fig:most_popular_words}
\end{figure}


\newpage

Among the 20 most popular words, the relative frequency in AI text ranges from 81\% to 124\% compared to their usage in human text. For example, the word "the" appears 81\% as often in AI text compared to human text. On the other hand, the word "as" is 124\% as frequent in AI text.

\vspace{3mm}

Apart from the most frequent words, some words are relatively overused in one subset compared to the other. These words are less frequent overall but still reveal stylistic patterns that can be utilized for classification.

\vspace{3mm}

Terms like "advertisement", "edit", and "percent" appear significantly more often in human-written texts, as shown in Figure~\ref{fig:words_human}. Also, humans use contractions like "that's" and "he's" (due to punctuation removal they are displayed as "thats" and "hes" respectively). AI models tend to avoid such contractions, using a more formal style instead.

\vspace{3mm}

\begin{figure}[h!]
    \centering
    \includegraphics[width=\textwidth]{Images/a_lot_human.png}
    \caption{Top 10 Words with Highest Relative Usage in Human Texts Compared to AI-Generated Texts.}
    \label{fig:words_human}
\end{figure}

\vspace{3mm}

On the other hand, AI-generated texts use words like "author", "stated", and "due" more often, which can be seen in Figure~\ref{fig:words_ai}. The full list shows a preference for formal and scientific language, which is typical of AI-generated content.

\vspace{3mm}

\begin{figure}[h!]
    \centering
    \includegraphics[width=\textwidth]{Images/a_lot_ai.png}
    \caption{Top 10 Words with Highest Relative Usage in AI-Generated Texts Compared to Human Texts.}
    \label{fig:words_ai}
\end{figure}


\vspace{3mm}

These differences, while being significant, do not directly translate into exceptional performance metrics. Modern AI-generated text is produced by state-of-the-art LLMs. Unlike deterministic generative methods like greedy or beam search, which might overuse most common words such as "the", "a", "and", the LLMs that generated text in our dataset, use stochastic techniques like top-k and top-p sampling. They ensure more varied generated text, reducing over-reliance on frequent words.

\vspace{3mm}

In comparison, more sophisticated models such as CNN and BiLSTM can also leverage occurrences of such words but do this in a slightly different manner. Both methods use fastText word embeddings, where each word is represented as the same vector. CNN and BiLSTM are capable of extracting information for classification not only from individual word embeddings but also from local combinations of words.  However, capturing relationships in longer sequences can be challenging, as these patterns may be rare, even in large~datasets.

\newpage

In CNNs, convolutional layers with filters of varying sizes (from 2 to 5) are applied, as detailed in the CNN model creation section~(Section~\ref{CNN model creation}). These filters capture patterns across sentences of 2 to 5 consecutive words, effectively learning n-gram-like features. 

\vspace{3mm}

For instance, not frequently used nouns in the dataset may have similar or less meaningful word embeddings. Such nouns are often preceded by determiners like "the", "a", or "an", forming bigrams that certain filters in the convolutional layer can capture. Similarly, filters of length 3 to 5 can identify longer patterns or phrases, such as "in other words", "however, while we can", or "as an AI language model". Despite their strength in capturing local structures, created CNN may struggle with long-range dependencies, as its architecture is focused on local features only.

\vspace{3mm}

BiLSTMs overcome this limitation, as they process sequences bidirectionally, extracting insights from both preceding and succeeding word contexts. Relationships between words can be extracted in a broader context. For instance, BiLSTMs can spot differences in meaning for polysemous words like "bank" in "river bank" versus "financial bank" based on the surrounding words. Despite their ability to handle sequential data, BiLSTMs rely on the quality of the input embeddings to generalize nuanced meanings.

\vspace{3mm}

Both CNN and BiLSTM architectures benefit from the use of fastText embeddings, which enhance their ability to represent individual words and subword structures. These embeddings provide clear representations that improve the models' performance, particularly for handling rare words or out-of-vocabulary terms. However, the extraction mechanisms of CNN and BiLSTM operate independently of the embeddings themselves, with fastText only being enhancement rather than the core mechanism of information~extraction.

\vspace{3mm}

In contrast, our BERT takes a fundamentally different approach to context extraction. Instead of relying on complex defined architectures like CNN or BiLSTM, BERT uses dynamic, context-sensitive embeddings that adapt to the surrounding words. This enables BERT to extract information at a deeper level with more flexibility. For instance, the meaning of "bank" is dynamically determined based on its usage in "river bank" or "financial bank", without requiring additional architecture complexity. In this case, the extraction process is driven by the Transformer layers, which detect relationships across the entire text, capturing both local and global dependencies. Then, it is wrapped by a simple neural network layer that doesn't add much complexity.


% For example, a lot of different nouns that are not frequent in the database could have a similar, not meaningful word-embedding. These nouns are often times written as successors of words such as "the", "a", "an". Then, these bigrams can be extracted from convolution layer by some created filters. Similar case for other filter lengths, where 3 to 5 sets of words can be extracted. This can extract longer sets of words, like "in other words", "however, while we can", "as an AI language model", etc.  

% In comparison, more sophisticated models such as CNN and BiLSTM can also leverage occurrences of such words, but do this in a slighlty different way. Both of these methods are used with FastText word-embeddings, so each word has the same vector representation. Both CNN and BiLSTM have abilities to not only extract information for classification, not only from single instances of these word-embeddings, but they do this locally, for combinations of words.

% In CNN, we add filters in convolutional layer, with sizes from 2, to 5. As described in CNN model creation section\ref{CNN model creation}, these filters can extract sets of 2 to 5 consequitive words. 

% In BiLSTM, we extract preceding and succeeding contexts, so information for classification is also 


% In comparison, more sophisticated models such as CNN and BiLSTM utilize occurrences of such words in a slightly different manner. Both methods leverage FastText word embeddings, where each word is represented as the same vector. CNN and BiLSTM are capable of extracting information for classification not only from individual word embeddings but also from local combinations of words.

% In CNNs, convolutional layers with filters of varying sizes (e.g., 2 to 5) are applied, as detailed in the CNN model creation section (\ref{CNN model creation}). These filters capture patterns across sequences of 2 to 5 consecutive words, effectively learning n-gram-like features.

% In BiLSTMs, information is extracted from both preceding and succeeding contexts, allowing the model to incorporate a broader view of the text for classification. This bidirectional structure enables BiLSTMs to better capture dependencies between words within a sentence.


\section{General Results}


Throughout the training of all models, the F1 score is selected as the primary evaluation metric. The F1 score, which represents the harmonic mean of precision and recall, offers a more nuanced assessment of incorrectly classified cases compared to the accuracy metric. This is particularly important in real-world classification problems where imbalanced class distributions are common, making F1 score a more reliable measure of a model’s~performance. 

\vspace{3mm}

The F1 score is calculated as follows:

\vspace{3mm}

\[
\text{F1-score} = 2 \cdot \frac{\text{Precision} \cdot \text{Recall}}{\text{Precision} + \text{Recall}}
\]

\vspace{3mm}

In our specific case, class imbalance is not a significant issue since the dataset is perfectly balanced, consisting of an equal number of human-written paragraphs and AI-rephrased paragraphs. However, the F1 score remains the preferred metric as it ensures balanced evaluation of both false positives (FP) and false negatives (FN). This balance is critical, as we aim to avoid significant disparities in distinguishing between these two types of~errors.

\vspace{3mm}

\noindent The F1 Score, as well as accuracy, precision and recall can be observed in Table~\ref{tab:performance-metrics}.

\vspace{3mm}

\begin{table}[h!]
    \centering
    \begin{tabular}{@{}lcccc@{}}
        \toprule
        \textbf{Model}            & \textbf{F1 Score} & \textbf{Accuracy} & \textbf{Precision} & \textbf{Recall} \\ \midrule
        Logistic Regression       & 0.887             & 0.885             & 0.872              & 0.904           \\
        SVM                       & 0.862             & 0.864             & 0.878              & 0.846           \\
        CNN                       & 0.904             & 0.902             & 0.885              & 0.924           \\
        BiLSTM (RNN)              & 0.934             & 0.935             & \textbf{0.954}     & 0.914           \\
        BERT                      & \textbf{0.963}    & \textbf{0.963}    & 0.952              & \textbf{0.974}  \\ \bottomrule
    \end{tabular}
    \caption{Performance Metrics}
    \label{tab:performance-metrics}
\end{table}

\vspace{3mm}

The evaluation of all five models reveals significant differences in performance across all measured metrics. Among all, the fine-tuned BERT demonstrates superior performance, achieving the highest results in F1 Score (96.3\%), Accuracy (96.3\%), and Recall (97.4\%). Especially, it presents an extremely high recall score, making it even more useful in scenarios where the primary objective is the detect nearly all AI-generated texts. For instance, this would be critically important in applications detecting AI-generated texts in sensitive domains such as news articles or academic submissions, where false negatives could have severe implications. In case of positive evaluation, particular work could be analyzed more deeply with a help of a human specialist.

\vspace{3mm}

Second best metrics are achieved by BiLSTM (RNN) model, with F1 Score of 93.4\%, Accuracy of 93.5\%, and the highest Precision (95.4\%). This model is particularly effective at minimizing the number of false positives in the dataset, however, it comes with the expense of a higher false-negative rate. This performance contrasts with that of BERT, as BiLSTM is particularly effective when avoiding the misclassification of human-written texts as AI-generated is a priority. For example, in cases like filtering emails for spam, where marking an important human-written message as spam could lead to it being lost, BiLSTM’s focus on precision makes it a better choice. While it does not perform as well overall as BERT, its ability to minimize false positives makes it valuable in situations where avoiding errors is more important than catching every AI-generated text.

\vspace{3mm}

It is important to note that in such specific cases, these models could be further optimized by training them with a customized loss function that assigns greater penalties to false negatives or false positives, depending on the context. Given that BiLSTM's precision is just 0.2\% higher than the one in BERT, the latter remains preferable choice.

\vspace{3mm}

With a similar basis to the BiLSTM model, including the use of fastText word embeddings, the CNN model also demonstrated rather strong performance, achieving an F1 Score of 90.4\% and an Accuracy of 90.2\%. However, it falls significantly behind BiLSTM by over 3\% in these metrics.

\vspace{3mm}

Traditional machine learning models such as Logistic Regression and SVM achieved the lowest performance among the tested models. Logistic Regression achieved an F1 Score of 88.7\% and Accuracy of 88.5\%, while SVM obtained an F1 Score of 86.2\% and Accuracy of 86.4\%. Nevertheless, these results are still impressive, falling only slightly behind the more complex deep learning model, CNN. Despite their inability to extract contextual information, these models leverage statistical features to make fairly accurate predictions. They might be particularly useful for creating lightweight models in a fast and simple~way.

\vspace{3mm}

An interesting observation is that SVM demonstrates an F1 Score 2.5\% lower than Logistic Regression, even though its process of model creation is generally considered more advanced. This highlights that increased model complexity does not always translate to better performance. In some cases, simpler models can outperform more complex ones, particularly when they are well-suited to the problem.


\section{Results Based on Text Length}

As highlighted in the data exploration section~(Section~\ref{Data Exploration}), input data for investigated labels varied and it is important to detect whether analyzed models worked correctly for varying text lengths. In order to do this, we group the test dataset into bins with range of 100 words and we compute accuracy on these subsets achieved by each model. Accuracy is chosen here because F1 score does not have this big utility here because most of the bins have big class imbalanced. The results are displayed in Figure
\ref{fig:word_count_vs_predictions}.

\vspace{3mm}

\begin{figure}[h!]
    \centering
    \includegraphics[width=\textwidth]{Images/word_count_vs_predictions.png}
    \caption{Average Accuracy of Models Across Token Length Bins}
    \label{fig:word_count_vs_predictions}
\end{figure}

\vspace{3mm}

The analyzed models exhibit similar accuracy trends across most text length ranges. Generally, model accuracy improves as text length increases, probably because longer texts provide more context and features for classification. For entries exceeding 600 words, all models achieve an accuracy of over 95\%. However, for simpler models like Logistic Regression and Support Vector Machines, this high accuracy is likely influenced by class imbalance, as these models do not effectively capture contextual information. Interestingly, in certain bins, these simpler models outperform more advanced ones, including BERT.

\vspace{3mm}

For the shortest texts, particularly those under 100 words, most models perform well, except for SVM, which struggles in this range. These strong performances are likely caused by the fact that the first bin only contains AI-generated entries, allowing models to leverage text length as a key feature.

\vspace{3mm}

In contrast, the bins containing texts between 100 and 300 words show the weakest performance. This range includes both AI- and human-generated entries, making it difficult for models to distinguish between sources based solely on length. Additionally, the shorter input length limits the models' ability to perform meaningful syntactic or contextual analysis, further reducing accuracy.

\vspace{3mm}

Lastly, it is worth noting that BERT's accuracy slightly decreases for very long texts. This suggests that BERT does not depend heavily on text length, which makes it more robust when handling long inputs, even if the training dataset is slightly biased. This resilience can be advantageous for real-world use cases involving diverse text lengths.

%%%%%

% As we can see, models present very similar dependencies between each other accross all text lengths. Models generally improve in accuracy as text length increases, which is likely due to longer texts providing more context and features for classification. For entries above 600 words, all models reached accuracy over 95\%. However, we have to keep in mind that this accuracy in the models that don't extract context well, that being Logistic Regression and Support Vector Machines, is probably caused mostly by class imbalance. Interestingly, for some bins they overpass these more complex models, even BERT. We can see that BERT's accuracy actually decreases for these long inputs, indicating that the model doesn't rely on text length and any other long input outside of the testing set could be tested without big bias. For shorter entries, we can see that in the first bin, for texts of less than 100 words, accuracy is very high in all models except for SVM. It is also probably caused by the fact that these entries are solely from AI source, so models rely in big proportion on the text length. In the latter bins we see that the models generally present the worst results here - these bins are well balanced and also the text are quite short so it's very hard to be accurate with predictions.

% We can focus on two best models - BERT and BiLSTM. BilSTM overally performes worse, however for texts over 600 characters it is superior, reaching almost perfect accuracy.  

% \section*{Performance Analysis and Conclusions}

% % XXXXX

% The evaluation of various models on the task of distinguishing AI-generated text from human-written text revealed significant differences in performance across metrics such as F1 Score, Accuracy, Precision, and Recall. Among all models, \textbf{BERT} demonstrated superior performance, achieving the highest scores in F1 Score (96.3\%), Accuracy (96.3\%), and Recall (97.4\%). These results highlight BERT's robust ability to effectively classify text with minimal errors.

% The \textbf{BiLSTM (RNN)} model ranked second, with a commendable F1 Score of 93.4\%, Accuracy of 93.5\%, and the highest Precision (95.4\%). This indicates that BiLSTM is particularly effective at minimizing false positives but slightly less capable of achieving the same recall levels as BERT. The \textbf{CNN} model also performed well, achieving an F1 Score of 90.4\% and Accuracy of 90.2\%, though it lags behind the performance of BERT and BiLSTM, especially in terms of Recall (92.4\%).

% In contrast, traditional machine learning models such as \textbf{Logistic Regression} and \textbf{SVM} demonstrated relatively lower performance. Logistic Regression achieved an F1 Score of 88.7\% and Accuracy of 88.5\%, while SVM obtained an F1 Score of 86.2\% and Accuracy of 86.4\%. These results suggest that traditional models, while effective to some degree, are less capable of capturing complex patterns in text compared to deep learning-based approaches.

% In summary, deep learning models, particularly BERT and BiLSTM, showcase superior performance in distinguishing AI-generated text from human-written text, with BERT leading across most metrics. These results emphasize the importance of leveraging advanced architectures like transformers for text classification tasks.

% % XXXXX



% \begin{figure}[h!]
%     \centering
%     \includegraphics[width=\textwidth]{Images/all_results.png}
%     \caption{Performance Metrics Comparison Across Models. This chart illustrates F1 Score, Accuracy, Precision, and Recall for each model.}
%     \label{fig:metrics-comparison}
% \end{figure}
% Let's break it down from the worst performing models to the best ones.



% Logistic regression:

% \begin{table}[h!]
% \centering
% \caption{Logistic Regression Metrics}
% \label{tab:performance_metrics}
% \begin{tabular}{|l|c|}
% \hline
% \textbf{Metric} & \textbf{Score} \\ \hline
% Accuracy Score & 0.8504 \\ \hline
% Confusion Matrix & 
% \begin{tabular}{c|c|c}
%      & Predicted Negative & Predicted Positive \\ \hline
%     Actual Negative & 5748 & 1228 \\ 
%     Actual Positive & 859 & 6117 \\ 
% \end{tabular} \\ \hline
% ROC AUC Score & 0.8504 \\ \hline
% F1 Score & 0.8543 \\ \hline
% \end{tabular}
% \end{table}

% CNN -- 0.90
% BiLSTM -- 0.94
% BERT -- 0.96